\section{Time Complexity}
\label{sec:time-complexity}

In this section, we analyze the \revisionrewrite{algorithms' time complexity}.

\textbf{\split.}
Conceptually, \split only performs both \revisionrewrite{optimization steps} exactly \emph{once}.  The worst-case complexity of algebraic optimization is covered by a large body of related work and determined to be exponential in the number of base relations~$n$~\cite{complexity_join_enumeration, np_hardness_join_enumeration}, even using state-of-the-art enumeration algorithms like \DPccp.  This can be highlighted by the fact that one needs to enumerate all subproblems, \ie elements of the powerset of all base relations. Furthermore, there are $\binom{n}{m}$ subproblems of size~$m$ and for each of them all possible splits have to be enumerated.  This can be simplified by choosing all subsets of size~1 up to size~$m-1$, \ie $\sum_{i=1}^{m-1}\binom{m}{i} = 2^m-2$. In total, we need to enumerate $\sum_{m=1}^{n}\binom{n}{m}*(2^m-2) = -2^{n+1}+3^n+1$ possible join orders, which resides in $\Bigo(3^n)$.

The following physical optimization step involves testing a finite number of physical operators for each algebraic operator contained in the received algebraic plan.  Moreover, a total order has to be determined using non-commutative cost models.  In the worst case, different sibling orderings result in different post-conditions that propagate throughout the entire optimization process.  For example, a physical join operator, receiving three or rather four QEPs with different post-conditions for its left and right subproblems, may yield up to 12 QEPs with different post-conditions.  However, the number of enumerated QEPs is still exponentially bounded by the number of joins as they are the only non-unary operator.  The number of joins is bounded by the number of base relations~$n$ minus~1.  Note that such an exponential propagation highly depends on the actual physical implementation, \eg a SHJ yields the post-condition of its probe child and thus stops the exponential propagation locally.  As both the algebraic and the physical optimization step require exponential time, \split's overall time complexity resides in~$\Bigo(3^n)$.

\textbf{\holistic.}
On the other hand, \holistic has to perform one conceptual physical optimization step \emph{per} found algebraic plan, \ie the exponential effort for physical optimization has to be spent multiple times instead of just once.  Even if implemented efficiently using dynamic programming, still all algebraic plan structures determining the join order together with all possible base relation permutations to achieve the total order, have to be enumerated.  The number of plan structures is asymptotically bounded by the Catalan number~\cite{book_moerkotte}, \ie exponentially, whereas the number of permutations is computed in~$\Bigo(n!)$ for worst-case query graph shapes.  Even if this complexity still resides in the exponential time class due to $n! \leq n^n = 2^{n*log(n)}$, the faculty renders \holistic to grow much faster than \split \wrt optimization time.  Our experimental evaluation in \autoref{sec:eval-opt_time} \revisionrewrite{confirms} our findings.

\section{Optimality}
\label{sec:optimality}

In this section, we analyze the \revisionrewrite{algorithms' optimality}.

\textbf{\split.}
The algorithm \split might get stuck in a local optimum due to its greediness of determining an algebraic plan without considering physical operators.  This greediness issue is already showcased by our running example in \autoref{sec:example}.  The algebraic optimization step decides to perform the join between~R and~S first because it results in the smaller intermediate result (see algebraic plan of \autoref{fig:example-plan-split}).  The succeeding physical optimization step cannot change the join order anymore and is only able to choose concrete physical operators (see QEP of \autoref{fig:example-plan-split}).  However, including the assumption from our running example regarding sortedness of~S.id and index on~T.sid, we have already shown that utilizing a MJ and a different join order yields a cheaper QEP \wrt the physical cost model (see \autoref{fig:example-plan-holistic}) that \holistic is able to find.

\textbf{\holistic.}
In contrast, \holistic enumerates all algebraic plan structures while testing all possible physical operators for each structure.  For our running example, this enumeration is depicted in \autoref{tbl:example-opt-holistic} and represents a superset of the enumeration of \split in \autoref{tbl:example-opt-split-phys}.  Therefore, \holistic is guaranteed to find the globally optimal QEP.  Even the recursion for nested query graphs does not prevent \holistic from being globally optimal \revisionadd{under the practical assumptions made} at the end of \autoref{sec:impl-fused_phys_ops}.

